% Generated by run_pdf_combo_sheet.py; do not edit.
\documentclass[10pt]{article}
\usepackage[table]{xcolor}
\usepackage{fontspec}
\usepackage{longtable}
\usepackage{array}
\usepackage{geometry}
\geometry{a4paper, landscape, left=1.2cm, right=1.2cm, top=1.3cm, bottom=1.3cm}
\setlength{\parindent}{0pt}
\definecolor{patchred}{RGB}{215,40,40}
\setmainfont{Libertinus Serif}[Path = ../fonts/, UprightFont = LibertinusSerif-Regular.otf, ItalicFont = LibertinusSerif-Italic.otf, BoldFont = LibertinusSerif-Semibold.otf, BoldItalicFont = LibertinusSerif-SemiboldItalic.otf]
\newfontfamily\Orig[Path = ../fonts/, UprightFont = LibertinusSerif-Regular.otf, ItalicFont = LibertinusSerif-Italic.otf, BoldFont = LibertinusSerif-Semibold.otf, BoldItalicFont = LibertinusSerif-SemiboldItalic.otf]{Libertinus Serif}
\newfontfamily\Patch[Path = ../fonts/, UprightFont = LibertinusSerif-Regular-patch.otf, ItalicFont = LibertinusSerif-Italic-patch.otf, BoldFont = LibertinusSerif-Semibold-patch.otf, BoldItalicFont = LibertinusSerif-SemiboldItalic-patch.otf]{Libertinus Serif Patch}
\newcommand\sheetbig{\fontsize{17}{22}\selectfont\rule[-10pt]{0pt}{33pt}}
\begin{document}
\section*{Test sheet: combining-mark combos in proofreading-633.pdf, original vs patched}
The PDF has 14 base\,+\,mark sequences in Libertinus Serif (7 distinct; none in semibold italic).
Every one was shaped with HarfBuzz in the original and the patched font of its style. The original reproduces the mark position printed in the PDF except in 1 combo, marked †, all in semibold: the PDF was set with another release of that font, so its ``original'' column is 7.051, not what the book printed. In semibold the PDF has the dot above on u (Sȯu̇th) 31 units higher than 7.051 puts it.
A combo is listed if the original font has no anchor for it (the mark is not anchored and is drawn after the base's advance) or the patch moves a mark by at least 10 units.
Each is shown in the first word of the PDF that has it, every character in the style the PDF set it in. Page = physical page of the PDF (1--420), not the book's pagination.
Columns: original, patched, and the two overlaid (patched in \textcolor{patchred}{red} over the original in black; only marks that move show red). Move = patched minus original mark position, (x, y) in font units, per mark.

\renewcommand\arraystretch{1}
\begin{longtable}{l r >{\sheetbig}l >{\sheetbig}l >{\sheetbig}l l >{\raggedright\arraybackslash}p{6.3cm}}
combo & page & \normalsize original & \normalsize patched & \normalsize overlay & move & occurrences \\ \hline \endhead
\multicolumn{7}{l}{\textbf{Regular}} \\
{\Orig\upshape\mdseries\char"0075\char"0307} {\tiny U+0075 U+0307} & 237 & {\Orig{\upshape\mdseries\char"0103\char"0063\char"0063\char"022F\char"0075\char"0307\char"02B9\char"006E\char"0074}} & {\Patch{\upshape\mdseries\char"0103\char"0063\char"0063\char"022F\char"0075\char"0307\char"02B9\char"006E\char"0074}} & \rlap{\color{patchred}{\Patch{\upshape\mdseries\char"0103\char"0063\char"0063\char"022F\char"0075\char"0307\char"02B9\char"006E\char"0074}}}{\Orig{\upshape\mdseries\char"0103\char"0063\char"0063\char"022F\char"0075\char"0307\char"02B9\char"006E\char"0074}} & {\footnotesize (+0,\,+31)} & {\footnotesize 5× (pp.~237); anchored} \\
{\Orig\upshape\mdseries\char"0075\char"0364} {\tiny U+0075 U+0364} & 314 & {\Orig{\upshape\mdseries\char"006E\char"006F\char"0074\char"0074\char"0075\char"0364\char"0072\char"0066\char"0066\char"0074\char"0069\char"0067}} & {\Patch{\upshape\mdseries\char"006E\char"006F\char"0074\char"0074\char"0075\char"0364\char"0072\char"0066\char"0066\char"0074\char"0069\char"0067}} & \rlap{\color{patchred}{\Patch{\upshape\mdseries\char"006E\char"006F\char"0074\char"0074\char"0075\char"0364\char"0072\char"0066\char"0066\char"0074\char"0069\char"0067}}}{\Orig{\upshape\mdseries\char"006E\char"006F\char"0074\char"0074\char"0075\char"0364\char"0072\char"0066\char"0066\char"0074\char"0069\char"0067}} & {\footnotesize (+16,\,+0)} & {\footnotesize 2× (pp.~314); anchored} \\
{\Orig\upshape\mdseries\char"006D\char"0303} {\tiny U+006D U+0303} & 53 & {\Orig{\upshape\mdseries\char"006A\char"00E4\char"006D\char"0303}} & {\Patch{\upshape\mdseries\char"006A\char"00E4\char"006D\char"0303}} & \rlap{\color{patchred}{\Patch{\upshape\mdseries\char"006A\char"00E4\char"006D\char"0303}}}{\Orig{\upshape\mdseries\char"006A\char"00E4\char"006D\char"0303}} & {\footnotesize (+12,\,+0)} & {\footnotesize 2× (pp.~53, 59); anchored} \\
\multicolumn{7}{l}{\textbf{Italic}} \\
{\Orig\itshape\mdseries\char"0075\char"0364} {\tiny U+0075 U+0364} & 259 & {\Orig{\itshape\mdseries\char"0062\char"0065\char"0079\char"0067\char"0065\char"0066\char"0075\char"0364\char"0067\char"0074\char"0065\char"005B\char"0072\char"005D}} & {\Patch{\itshape\mdseries\char"0062\char"0065\char"0079\char"0067\char"0065\char"0066\char"0075\char"0364\char"0067\char"0074\char"0065\char"005B\char"0072\char"005D}} & \rlap{\color{patchred}{\Patch{\itshape\mdseries\char"0062\char"0065\char"0079\char"0067\char"0065\char"0066\char"0075\char"0364\char"0067\char"0074\char"0065\char"005B\char"0072\char"005D}}}{\Orig{\itshape\mdseries\char"0062\char"0065\char"0079\char"0067\char"0065\char"0066\char"0075\char"0364\char"0067\char"0074\char"0065\char"005B\char"0072\char"005D}} & {\footnotesize (+15,\,-1)} & {\footnotesize 1× (pp.~259); anchored} \\
\multicolumn{7}{l}{\textbf{Semibold}} \\
{\Orig\upshape\bfseries\char"0075\char"0307} {\tiny U+0075 U+0307} & 231 \textsuperscript{†} & {\Orig{\itshape\mdseries\char"0053}{\upshape\bfseries\char"022F\char"0075\char"0307}{\itshape\mdseries\char"0074\char"0068}} & {\Patch{\itshape\mdseries\char"0053}{\upshape\bfseries\char"022F\char"0075\char"0307}{\itshape\mdseries\char"0074\char"0068}} & \rlap{\color{patchred}{\Patch{\itshape\mdseries\char"0053}{\upshape\bfseries\char"022F\char"0075\char"0307}{\itshape\mdseries\char"0074\char"0068}}}{\Orig{\itshape\mdseries\char"0053}{\upshape\bfseries\char"022F\char"0075\char"0307}{\itshape\mdseries\char"0074\char"0068}} & {\footnotesize (-1,\,+37)} & {\footnotesize 1× (pp.~231); anchored} \\
\end{longtable}
\section*{Not covered by the patch: reproductions}
\subsection*{p.~263}A second above mark on a precomposed letter that already has one. Italic ö and ä have no above anchor in either font (stage 9 inherits an above anchor only for letters without an above mark), so the grave is not anchored and lands on the diaeresis; regular ä has a stacking anchor, and so does italic ā (ā́ in 614.pdf is anchored). Not covered by the patch. Here in italic \emph{die Jä̀gerey}, against the regular ä̀ after it.\par\smallskip
\begin{tabular}{l >{\sheetbig}l >{\sheetbig}l >{\sheetbig}l}variant & \normalsize original & \normalsize patched & \normalsize overlay \\ \hline
as printed & {\Orig{\itshape\mdseries\char"004A\char"00E4\char"0300\char"0067\char"0065\char"0072\char"0065\char"0079}} & {\Patch{\itshape\mdseries\char"004A\char"00E4\char"0300\char"0067\char"0065\char"0072\char"0065\char"0079}} & \rlap{\color{patchred}{\Patch{\itshape\mdseries\char"004A\char"00E4\char"0300\char"0067\char"0065\char"0072\char"0065\char"0079}}}{\Orig{\itshape\mdseries\char"004A\char"00E4\char"0300\char"0067\char"0065\char"0072\char"0065\char"0079}} \\
regular ä̀ & {\Orig{\upshape\mdseries\char"00E4\char"0300}} & {\Patch{\upshape\mdseries\char"00E4\char"0300}} & \rlap{\color{patchred}{\Patch{\upshape\mdseries\char"00E4\char"0300}}}{\Orig{\upshape\mdseries\char"00E4\char"0300}} \\
\end{tabular}\par\medskip
\subsection*{Checked in context: not font problems}p.~231, Table~7, \emph{Sȯu̇th}: italic S and th, with ȯ (precomposed U+022F) and u + U+0307 in upright semibold (Buchanan's digraph diacritics, set bold on purpose). In the PDF the two dots are level: the PDF's dot on u is 31 units above where 7.051 puts it (see above), and the patch raises it 37, so the patched font also sets them level. Not a problem in the PDF.\par\medskip
\section*{Other combos (moved under 10 units, not set)}\small
Regular {\Orig\upshape\mdseries\char"00E4\char"0300} p.~263 (1×).
\end{document}
